\section{Introduction}

Every wasted FLOP has a dollar cost that compounds at datacenter scale. Anthropic is scaling from 2.5~GW to over 5~GW by year-end at \$10--13B/GW/year; the Big Four hyperscalers are spending a combined \$600B in CapEx this year. The problems are coupled: generation throughput gates RL rollouts, algorithm choice determines whether those rollouts produce signal, and the unit economics---cost per useful training step, hardware utilization, the 5--6 year payback per GW deployed---determine whether a research direction is financially viable.

Anthropic's RL infrastructure spans performance engineering, research engineering, and research science. Understanding across these boundaries makes you better at any one of them.

This guide is organized around a single metric: \textit{intelligence per Watt}. Every chapter builds cross-boundary intuition and produces a publishable artifact.

% --- Governing Equation ---

\begin{tcolorbox}[colback=matcha!5!white, colframe=matcha!70!black, boxrule=0.6pt,
  left=5pt, right=5pt, top=6pt, bottom=6pt,
  fonttitle=\bfseries, title=The Governing Metric]
\[
\underbrace{\frac{\text{Intelligence}}{\text{Watt}}}_{\text{What we maximize}}
\;=\;
\underbrace{\left(\frac{\text{Quality}}{\text{FLOP}}\right)}_{\substack{\text{Ch\,5--6:}\\\text{RL Depth}}}
\;\times\;
\underbrace{\left(\frac{\text{FLOPs}}{\text{Watt}}\right)}_{\substack{\text{Ch\,1--4,\,7:}\\\text{Hardware \& Systems}}}
\]
\smallskip
The ``per Watt'' is both physics and dollars. At \$10--13B per GW per year, a 10\% improvement in utilization is worth hundreds of millions.
\end{tcolorbox}

% --- Three-Tier Hierarchy ---

\subsection{How the Chapters Compose}

The chapters do not compose sequentially --- they compose \textit{hierarchically}. The hardware chapters set the iteration speed (experiments per week). The algorithm chapters determine what each experiment is testing. Both must be healthy for a research program to move fast.

\begin{center}
\begin{tikzpicture}[
  tier/.style={draw, thick, rounded corners=4pt, minimum height=1.4cm, text centered, font=\small},
  arrow/.style={-{Stealth[length=6pt]}, thick, sumi!60}
]
  % Tier 1 — base
  \node[tier, fill=aiiro!8, minimum width=12cm] (t1)
    {\textbf{Hardware \& Systems} (Ch\,1--4,\,7) --- FLOPs/Watt};

  % Tier 2 — top
  \node[tier, fill=beni!8, minimum width=8cm, above=0.6cm of t1] (t2)
    {\textbf{RL Depth} (Ch\,5--6) --- Quality/FLOP};

  % Arrows
  \draw[arrow] (t1.north) +(-2,0) -- +(-2,0.6);
  \draw[arrow] (t1.north) +(2,0) -- +(2,0.6);

  % Label
  \node[font=\footnotesize\itshape, sumi!70, right=0.3cm of t2] {builds on};
\end{tikzpicture}
\end{center}

% --- PE-RL Flywheel ---

\subsection{The PE--RL Flywheel}

Performance engineering and reinforcement learning are not separate concerns. They form a reinforcing loop:

\begin{center}
\begin{tikzpicture}[
  node distance=5cm,
  box/.style={draw, thick, rounded corners=6pt, fill=#1, minimum width=3.2cm, minimum height=1.6cm, text centered, font=\small\bfseries, text width=2.8cm},
  lbl/.style={font=\footnotesize, text width=3.2cm, align=center}
]
  \node[box=aiiro!10] (pe) {Performance\\Engineering};
  \node[box=beni!10, right=of pe] (rl) {Reinforcement\\Learning};

  \draw[-{Stealth[length=7pt]}, thick, bend left=25]
    (pe.north east) to node[lbl, above] {Higher throughput\\$\rightarrow$ more rollouts/hr} (rl.north west);
  \draw[-{Stealth[length=7pt]}, thick, bend left=25]
    (rl.south west) to node[lbl, below] {Policy learns kernel\\configs, quantization} (pe.south east);
\end{tikzpicture}
\end{center}

\noindent Higher generation throughput means more gradient steps per week, which means faster policy improvement --- or faster discovery that a hypothesis is wrong. Conversely, a policy trained with verifiable rewards on compiler or hardware benchmarks can learn optimization decisions that outperform human heuristics. The person who understands both sides of this loop is uniquely valuable.

% --- Chapter Map ---

\subsection{Chapter Map}

\begin{enumerate}
  \item \textbf{CUDA Kernel Optimization.} Foundational GPU kernel engineering on NVIDIA hardware. The prerequisite for every other hardware chapter. Produces a Boehm-style kernel progression with roofline analysis.

  \item \textbf{Benchmarking AMD MI300X.} Evaluates a multi-billion dollar hardware investment: where does AMD win on intelligence-per-Watt, and what engineering investment closes the gaps? Produces a decision-grade benchmarking report.

  \item \textbf{TPU/Pallas Kernel Optimization.} Google's answer to the same efficiency question. Covers Anthropic's primary accelerator fleet and contrasts with the GPU approach of Chapters~1--2.

  \item \textbf{RL Inference \& Serving Efficiency.} The \#1 RLVR bottleneck: autoregressive generation consumes 50--70\% of training wall time. Hardware-agnostic techniques for the throughput problem that gates everything downstream.

  \item \textbf{Policy Gradients --- Hypothesis-Driven Depth.} Which policy gradient choices --- clipping, GAE, KL penalties, entropy regularization --- are actually load-bearing, and why? Produces confirmed or refuted hypotheses with clean ablations. \textit{Prerequisites: none. Can be started in parallel with Chapter~1 at \$0 compute cost.}

  \item \textbf{TinyRL --- Community OSS Framework.} Chapter~5 theory operationalized in a community-auditable framework. Every design decision has a clean ablation. Produces an open-source artifact with leaderboard entries.

  \item \textbf{Distributed Training \& Communication.} The capstone: scaling everything from Chapters~1--6 across thousands of nodes. Communication overhead, parallelism strategies, and cluster utilization at 5+~GW scale.
\end{enumerate}

% --- Portfolio Comparison ---

\subsection{Choosing Your Portfolio Project}

\begin{table}[h]
\centering
\footnotesize
\setlength{\tabcolsep}{3pt}
\begin{tabular}{@{}lccccccc@{}}
\toprule
\textbf{Criterion} & \textbf{CUDA} & \textbf{AMD} & \textbf{Pallas} & \textbf{Infer.} & \textbf{PG} & \textbf{TinyRL} & \textbf{Distrib.} \\
 & \textbf{(Ch\,1)} & \textbf{(Ch\,2)} & \textbf{(Ch\,3)} & \textbf{(Ch\,4)} & \textbf{(Ch\,5)} & \textbf{(Ch\,6)} & \textbf{(Ch\,7)} \\
\midrule
Clean feedback loop?   & Yes  & Yes   & Yes    & Yes  & Yes     & Yes  & Yes  \\
Fast iteration?        & Yes  & Yes   & Med.   & Yes  & Yes     & Yes  & Med. \\
Underexplored?         & Low  & V.Hi  & High   & High & ?       & High & Med. \\
RLVR-relevant?         & Yes  & Yes   & Yes    & Very & Yes     & Very & Yes  \\
\bottomrule
\end{tabular}
\end{table}

\noindent\textbf{By goal:}
\begin{itemize}
  \item \textbf{Fastest artifact (5--7 weeks):} Ch\,2 AMD benchmark. One number --- MI300X vs.\ H100 on end-to-end RLVR throughput --- starts a conversation.
  \item \textbf{Strongest community signal:} Ch\,6 TinyRL. Open-source frameworks with leaderboard entries are persistent artifacts that compound.
  \item \textbf{Broadest role targeting:} Ch\,1 + Ch\,3. Covers GPU and TPU kernel engineering simultaneously.
  \item \textbf{Highest research ceiling:} Ch\,5 policy gradient hypotheses. A confirmed or refuted hypothesis about clipping behavior or GAE bias is more valuable than a benchmark table, if it changes how people train policies.
\end{itemize}

% --- Methodology ---

\subsection{Methodology}

Every project in this guide follows the same template:

\begin{enumerate}
  \item Pick a well-defined problem with clean verification.
  \item Build the simplest version that could possibly produce a signal.
  \item Ablate one variable at a time. Predict before you observe.
  \item Document honestly --- negative results and failed hypotheses are publishable.
\end{enumerate}

% --- Box Types ---

\subsection{How to Read This Guide}

Chapters use three box types to distinguish different kinds of work:

\begin{thinkfirst}{Example}
\textit{Conceptual foundations.} Reason through the problem before touching code. No computer needed.
\end{thinkfirst}

\begin{exercise}{Example}
\textit{Hands-on implementation.} Explicit deliverables and verification criteria.
\end{exercise}

\begin{portfolio}{Example}
\textit{Publishable projects.} External verification criteria and a clear artifact (blog post, OSS repo, benchmark report).
\end{portfolio}
